
Dataset versioning tools provide infrastructure for tracking changes but lack semantic meaning—researchers cannot determine from version identifiers whether updates represent breaking changes requiring model retraining. We tested whether statistical drift detection using fixed thresholds derived from ImageNet literature could automate semantic version classification (MAJOR/MINOR/PATCH) for machine learning datasets. Experiments on 14 dataset pairs (1 vision, 13 NLP) evaluated a classifier using Kolmogorov-Smirnov test and Maximum Mean Discrepancy on PCA-reduced features with thresholds of 7\%, 2\%, and 0.5\% for MAJOR, MINOR, and PATCH classifications respectively. The classifier achieved 28.57\% overall accuracy against literature-derived ground truth labels, with 25\% precision and 25\% recall for MAJOR change detection. This performance falls below the 33\% random baseline for three classes and substantially below the predefined success criteria of 85\% accuracy, 70\% precision, and 85\% recall. Root cause analysis identified that drift scores varied by 18.8× across datasets (range 0.042 to 0.79), with PATCH-labeled datasets scoring higher than some MAJOR-labeled datasets. The classifier misclassified 3 out of 4 MAJOR changes and incorrectly flagged 6 out of 7 non-MAJOR changes as MAJOR. These results falsify the hypothesis that ImageNet-derived thresholds generalize across dataset types and suggest that drift magnitude is dataset-relative rather than absolute.

\textbf{Keywords:} dataset versioning, distribution shift detection, semantic versioning, reproducibility, negative result
